\chapter{LBM 的现代模型与算法}\label{ch:16}
\begin{chaptersummary}
单松弛等温 LBM 已给出离散速度到宏观输运的路线。本章通过矩、累积量和离散熵比较不同碰撞设计，再讨论热、可压缩与边界模型需要增加的自由度和矩条件。模型选择由目标物理、Mach 数、温度变化与可验证的稳定范围共同决定。
\end{chaptersummary}
\learninggoals{区分原始矩、中心矩和累积量；写出 MRT 与离散熵约束；解释热与可压缩模型的额外矩需求；设计边界和 Galilean 误差检验。}
\studyroute{给出模型结构和基本公式，略去具体 MRT 矩阵表、累积量高阶回变换和特定可压缩格子的系数；这些内容按所选原始模型展开。}{MRT lattice Boltzmann; cumulant collision; entropic lattice Boltzmann; thermal lattice; Galilean invariance}{比较 BGK、MRT 与累积量碰撞保持哪些矩、分别调节哪些非守恒模式，再为同一剪切波规定统一黏度与误差指标。}

\section{MRT：在矩空间中调节松弛}\index{MRT--LBM}
对可逆矩变换 $m=Tf$，多松弛时间碰撞为
\[
m^*=m-S(m-m^{\mathrm{eq}}),\qquad
f^*=T^{-1}m^*.
\]
$S=\diag(s_a)$ 的守恒矩对应零松弛率，其余模式可有不同速率。
与剪切应力对应的速率 $s_\nu$ 决定
\[
\nu=c_s^2\Delta t\left(\frac1{s_\nu}-\frac12\right).
\]
其他非守恒模式的调节影响声学、体积黏性和格子尺度阻尼。
对局部对角松弛，$0<s_a<2$ 控制单独松弛模态；
完整更新还含迁移、平衡态线性化与边界，需检查组合谱~\cite{dhumieres2002}。

\conceptbridge{为何多松弛模型要指定所用矩，而不能只列出一组松弛率？}
{矩变换决定哪些模式对应守恒、应力和更高阶非平衡信息。
同一组速率作用在不同矩上会产生不同碰撞算子。
基于中心矩或累积量的碰撞还改变了变量与局部速度的关系。}
{碰撞坐标 $\longrightarrow$ 模态松弛 $\longrightarrow$ 输运与非流体误差。}

\section{中心矩与累积量}
原始矩为 $\sum_i f_i\bm c_i^\alpha$，中心矩以
$(\bm c_i-\bm u)^\alpha$ 替代。
若 $\rho>0$，用归一化粒子群 $f_i/\rho$ 定义生成函数
\[
K(\bm t)=\log\left(\frac1\rho\sum_i f_i
\exp[\bm t\cdot(\bm c_i-\bm u)]\right),\qquad
\kappa_\alpha=\left.\partial_{\bm t}^{\alpha}K\right|_{\bm t=0}.
\]
二阶累积量给出中心协方差，高阶累积量扣除了低阶矩乘积的贡献。
累积量碰撞在这些坐标中松弛，再重建 $f_i$，
从而调节不同阶非平衡信息的耦合~\cite{geier2015}。
有限速度集合仍限制可独立表达的矩，各模型的平衡态、速率与修正项需要一起指定。

\section{离散熵与正性}\index{熵格子 Boltzmann}
对正粒子群，可定义凸函数
\[
H(f)=\sum_i f_i\log\frac{f_i}{w_i}.
\]
熵平衡态在质量和动量约束下最小化 $H$。
一种熵碰撞沿 $\Delta=f^{\mathrm{eq}}-f$ 寻找保持正性的非零根
$H(f+\alpha\Delta)=H(f)$，再用 $0<\beta\le1$ 取
$f^*=f+\alpha\beta\Delta$。
凸性给出该线段上的 $H(f^*)\le H(f)$。
上述熵估计以熵平衡态、正性域及非零根存在为前提。采用多项式平衡态或固定 BGK 速率时，应对相应碰撞映射重新建立这些条件。
原始离散熵构造见~\cite{karlin1999}。

\section{热模型、可压缩模型与边界}
等温格子的 $p=c_s^2\rho$ 固定了相应状态关系。
热传输需增加能量相关的离散分布或更丰富速度集合，
并匹配热流矩、比热和目标 Prandtl 数。
可压缩流还需正确的能量通量、可变温度响应及足够阶的速度矩；
具体模型按速度集合、平衡展开与碰撞形式验证其宏观极限。

静止壁面的半程反弹将出射粒子群反向为入射粒子群；
壁面位置、滑移误差与几何处理相关。移动壁面需补动量交换，
入口与出口则需按给定宏观条件重建未知入射群。
热边界还需给定能量交换。曲面插值、漫反射与稀薄气体边界有各自适用条件。

\section{稳定性与 Galilean 误差的比较}
对剪切波叠加均匀背景速度 $U_b$，连续流体的衰减率应与 $U_b$ 无关。
比较不同背景速度下的黏度与相位，可测量格子高阶矩造成的 Galilean 误差。
同时记录 Mach 数、黏度、格子分辨率、边界位置及粒子群正性。
稳定性分析可结合线性谱、离散熵与非线性实验，不同指标针对不同误差机制。

\aiexercise{从 D2Q9 中选择一个明确的可逆矩变换，标出守恒与剪切模式。令非守恒速率相同，核对其在相应平衡矩下何时恢复 BGK 更新。}
{对周期剪切波添加两个不同背景速度，比较 BGK、MRT 和一个有出处的累积量模型的衰减误差；再说明若要计算温度变化还缺哪些状态量与矩约束。}
